% Sect. 4 Pipeline (outline brief 6.4). Describes the tool at commit 5ffd56e of the main branch ("Work inside the
% DCAT state each DCAT-AP release builds on; keep all versions; rebase"), which is not merged into the paper branch
% (TODO in Sect. 4.4). In the src comments, B = backend/profile_extensioner at 5ffd56e; CONF5, DCS, ENG, ISS, KCU,
% LIT, LITD are the outline's shorthands (docs/conformance.md at 5ffd56e; scratchpad digests dcat-state, engine,
% issues-at-main, knowledge-compile-ui; docs/literature-steps.md; literature digest). Mirrors RQ1 Sect. 4
% (L236-355): intro, Fig. 1, three subsections that follow the rows of the figure, then the implementation.
\section{Pipeline}\label{sec:pipeline}

\begin{figure}[t]
  \centering
  % Fig. 1: the pipeline of the tool at commit 5ffd56e. Mirrors the RQ1 paper's Fig. 1: numbered operations in three
% labelled rows, the middle row read from right to left. Encoding in RQ1's grammar: dashed border = the operation
% calls the model (3 proposer, 4 judge, 5 repairer); light fill = a person decides (9); plain = program code. The
% blackboard band is program code as well, so it has a plain (solid, grey) border.
% Sources of the labels: B/engine/intake.py, retrieve.py, propose.py, verify.py (V1-V5, D1-D3, T1, B1),
% pipeline.py (T2 judge, repair, consolidation, decision, sign-off), compile/, conformance/check.py (C1-C8)
% at 5ffd56e; B = backend/profile_extensioner.
\begin{tikzpicture}[
    font=\scriptsize,
    op/.style={draw, rounded corners=2pt, minimum height=1.08cm, align=center,
               inner sep=1.5pt, text width=3.7cm},
    llm/.style={op, dashed},
    human/.style={op, fill=black!8},
    row/.style={font=\scriptsize\bfseries, anchor=south west, inner sep=1pt},
    flow/.style={-{Stealth[length=1.5mm]}, semithick},
    side/.style={-{Stealth[length=1.5mm]}, semithick, dashed},
  ]
  % row 1: read and propose (left to right)
  \node[op]  (o1) at (0,0)     {1.\ Intake: lint the atoms,\\ map subjects to classes};
  \node[op]  (o2) at (4.15,0)  {2.\ Retrieve candidates: the\\ base, its DCAT, other\\ vocabularies};
  \node[llm] (o3) at (8.3,0)   {3.\ Propose a disposition and\\ typed changes per atom; B1\\ checks covered-by-base claims};
  % row 2: check and repair (right to left)
  \node[llm] (o4) at (8.3,-1.75)  {4.\ Verify: V1--V5, D1--D3, T1\\ in code, T2 by a model judge};
  \node[llm] (o5) at (4.15,-1.75) {5.\ Repair failing changes,\\ at most three rounds}; % src: @5ffd56e B/contracts.py:595
  \node[op]  (o6) at (0,-1.75)    {6.\ Consolidate: one element\\ per class and property};
  % row 3: compile and release (left to right)
  \node[op]    (o7) at (0,-3.7)    {7.\ Compile SHACL, LinkML,\\ DSV, ProfileDelta, docs;\\ V6 checks the SHACL};
  \node[op]    (o8) at (4.15,-3.7) {8.\ Conformance C1--C8,\\ witness against the\\ official shapes};
  \node[human] (o9) at (8.3,-3.7)  {9.\ Review the queue, decide,\\ sign off the release};
  % row labels
  \node[row] at (o1.north west) {Read and propose};
  \node[row] at (o6.north west) {Check and repair};
  \node[row] at (o7.north west) {Compile and release};
  % flow
  \draw[flow] (o1) -- (o2);
  \draw[flow] (o2) -- (o3);
  \draw[flow] (o3) -- (o4);
  \draw[flow] ([yshift=1.5mm]o4.west) -- ([yshift=1.5mm]o5.east);
  \draw[side] ([yshift=-1.5mm]o5.east) -- ([yshift=-1.5mm]o4.west);
  \draw[flow] (o5) -- (o6);
  \draw[flow] ([xshift=13mm]o6.south) -- ([xshift=13mm]o7.north);
  \draw[flow] (o7) -- (o8);
  \draw[flow] (o8) -- (o9);
  \draw[side] (o4.south) -- node[left=1pt, font=\tiny, align=right] {queue: errors after repair, low confidence,\\
    conflicts, later-DCAT terms} (o9.north);
  % blackboard (program code: plain border)
  \node[draw=black!55, rounded corners=2pt, anchor=north, text width=11.75cm, align=center,
        inner sep=3pt] (bb) at (4.15,-4.45)
        {Blackboard: append-only log of every rule, model call and human action\\
         current state of every atom and change\ $\cdot$\ model calls with prompts, answers, tokens and a replay cache};
\end{tikzpicture}

  \caption{The pipeline. Dashed operations call the model; with the rules baseline, offline rules take their place
    (Sect.~\ref{sec:setup}). Light: the person's operation. A change that fails a check at error level goes to
    repair~(5), at most three rounds, and, still failing, to the review queue~(9), as do conflicts, changes below % src: @5ffd56e B/contracts.py:595 (repair_rounds 3); B/engine/pipeline.py:246-267, 295-338, 347-353
    the confidence threshold and, by default, changes that use a term of a later \acs{DCAT}. Proposer, repairer and judge use the
    same model, so the judge is not independent of the proposer.}
  \label{fig:pipeline}
\end{figure}

Figure~\ref{fig:pipeline} shows the nine operations in three rows: the model acts as proposer~(3), judge~(4) and % src: paper/figures/pipeline.tex (nodes o1-o9, three rows); @5ffd56e B/contracts.py:104-107 (nine stages); B/engine/prompts.py:1 (three roles)
repairer~(5), a person decides in the last~(9), and program code does the rest. The supplementary material follows % src: paper/figures/pipeline.tex (operation numbers 3, 4, 5, 9)
one atom of a synthetic demo through them.
\todo{supplementary material: one atom of the synthetic demo through the nine operations}
The pipeline also runs without a model, with offline rules in the dashed operations (Sect.~\ref{sec:setup}). % src: @5ffd56e B/llm/mock.py:1-10

\subsection{Read and propose}\label{sec:read}

Intake checks the form of every atom (Sect.~\ref{sec:problem}) without rewriting it and maps its subject to a % src: @5ffd56e B/engine/intake.py:1-7, 108-131
class of the base. An \ac{IRI} or a prefixed name in the statement is an error, and the atom stays unclear % src: @5ffd56e B/engine/intake.py:39-51, 115-116; B/engine/pipeline.py:184-185
without a proposal; a subject that is not a plain word, a very short statement, a compound statement and a % src: @5ffd56e B/engine/intake.py:52-87 (warnings; vague words and an unspecified modality are info)
modality the wording contradicts are warnings, the last two adapted from characteristics of a well-formed % LIT G2-S02 (partial, Ours): compound = singular, modality mismatch = unambiguous
requirement \cite{iso29148}.

Retrieval is deterministic: a field-weighted lexical score, without embeddings. The base's elements on the atom's % src: @5ffd56e B/knowledge/vocab_index.py:8, 332-343, 352-367 (label 3, local name 2, definition 1; idf); KCU §2.4
class are boosted and, when they match reasonably, listed first; properties that the base's \acs{DCAT} lists for % src: @5ffd56e B/engine/retrieve.py:40, 52, 108-132, 259-263 (BASE_FIRST_SHARE 0.35)
the class but the base does not profile are boosted ($\times$1.4; the 1.x bases have none); terms that only a later % src: @5ffd56e B/engine/retrieve.py:41, 134-156; DCS §2.4 (rung 2 empty for 1.x)
\acs{DCAT} defines are marked with their version and down-weighted ($\times$0.5), and terms the base deprecates % src: @5ffd56e B/engine/retrieve.py:42, 82-88; B/engine/prompts.py:136-138
and vocabularies made for one domain are down-weighted too. % src: @5ffd56e B/engine/retrieve.py:43, 47-49, 186-190

The proposer is called once per atom, without the other atoms' changes, and fills a typed record by a numbered % src: @5ffd56e B/engine/pipeline.py:188-191 (no pending changes passed); ISS E7.2
procedure: it restates the information needed, fixes the class and walks the ladder of E1 from the base slice to a % src: @5ffd56e B/engine/prompts.py:30-48
mint, taking a later-\acs{DCAT} term only when no other term carries the information and minting only after saying why each plausible
candidate fails; every constraint needs a reason, and the answer a confidence. Program code then builds the change
records of Sect.~\ref{sec:problem}, and a refinement inherits from the base what it leaves unset, except the % src: @5ffd56e B/engine/propose.py:175-306 (unchanged since 2b92d77, ISS header)
obligation. B1 checks each covered-by-base claim as E1 requires, and a failed claim leaves the atom unclear. % src: @5ffd56e B/engine/verify.py:513-533; B/engine/pipeline.py:200-211
One atom per memoryless call that first asks whether the base already meets it follows the released Memoryless % LIT G1-S03 (Ours: the merge), G1-S05
CQbyCQ prompt \cite{lippolis2025ontology}, and reuse limited to retrieved terms follows RIGOR's released % LIT G1-S08
implementation \cite{fathallah2026rigor}; the ladder is our reading of SEMIC's guidance \cite{semicreuse2026}. Here the model % LIT G3-S01 (implemented, Ours)
chooses the properties, within typed records, and never writes \ac{SHACL} or mints an identifier.

\subsection{Check and repair}\label{sec:check}

Program code checks every change in a fixed order (Table~\ref{tab:checks} in the appendix): V1 term existence, V2 % src: @5ffd56e B/engine/verify.py:3-18, 166-211, 482-483; ISS E7.1 (union domains fixed)
domain and range, union domains included, D1 to D3 the \acs{DCAT} state, V3 legal refinement, V4 reuse before mint,
V5 consistency and T1 the trace; a set-level V5 error marks every live change on the same class and property. % src: @5ffd56e B/engine/verify.py:493-510
By default a term that only a later \acs{DCAT} defines gives a D1 warning that sends the change to a person; a % src: @5ffd56e B/contracts.py:604-606 (later_dcat_terms review|forbid|allow); B/engine/verify.py:228-242
term only an earlier \acs{DCAT} defines and a \acs{DCAT} property without a declared domain on a class the base's % src: @5ffd56e B/engine/verify.py:243-245, 259-273
\acs{DCAT} does not list it for are D1 errors, terms the base's \acs{DCAT} or the base deprecates are D2 errors, % src: @5ffd56e B/engine/verify.py:246-257
and D3 warns on every use of the six inverse properties in the \acs{DCAT} namespace. Only a change without an % src: @5ffd56e B/engine/verify.py:274-280; DCS §2.6 (11 DCAT 3 pairs, 5 outside the dcat: namespace)
error reaches the judge T2, which is told to judge from the atoms as written, not from published profiles: % src: @5ffd56e B/engine/pipeline.py:279-282; B/engine/prompts.py:84-94
supported passes, unsupported or a named unsupported constraint is an error, and partially supported or an
unavailable judge is a warning. We adapt checking before judging from the released code of a recent blackboard orchestration % src: @5ffd56e B/engine/pipeline.py:221-244; LIT G4-S04 (Ours)
of \acp{LLM} \cite{chmielewski2026blackboard}\verify{chmielewski2026blackboard: known from its released code only;
the entry has no DOI or URL}, and make a missing trace an error where RIGOR's released implementation warns % LIT G1-S18 (Ours)
\cite{fathallah2026rigor}.

A change with an error goes to repair, at most three rounds, scoped to its failed checks as in the released code % src: @5ffd56e B/contracts.py:595; B/engine/pipeline.py:295-338; B/engine/prompts.py:238-242
of RIGOR and IDEA2 \cite{fathallah2026rigor,alharbi2027idea2}: the repairer revises % LIT G1-S13; G2-S13
it, and the new version is verified, judged and decided again, or withdraws it. Consolidation then merges % src: @5ffd56e B/engine/pipeline.py:311-333
compatible changes on one class and property, with the strongest obligation and the tightest bounds, and queues
incompatible ones as a conflict. A change is auto-accepted when no check failed at error level, automatic % src: @5ffd56e B/engine/consolidate.py:25-95 (unchanged since 2b92d77); B/engine/pipeline.py:340-366
acceptance is on, its self-reported confidence is at least 0.5, a majority agrees if there are several votes, and % src: @5ffd56e B/engine/pipeline.py:246-267; B/contracts.py:598-599
no D1 warning asks for a person; otherwise it waits for review. Other warnings never block, and the threshold is % LITD §7 (no cited work supports self-reported confidence as calibrated)
our setting.
\todo{tool: state or fix: duplicates on one class and property go to repair before consolidation would merge them;
a change cleared by consolidation is decided without a judge verdict on its merged atoms (ENG §7.2-7.3; ISS E7.2,
E7.3; Q22)}

\subsection{Compile and release}\label{sec:compile}

Only included changes compile, into seven artifacts. The \ac{SHACL} delta holds only what the extension adds or % src: @5ffd56e B/compile/model.py:47-49; B/engine/pipeline.py:382-392 (seven artefacts)
tightens, with a separate \term{sh:minCount}~1 shape at Warning for a recommended property and, on each property shape, the reason and % src: @5ffd56e B/compile/shacl.py:3-8, 40-91
a \term{prov:wasDerivedFrom} link per covered atom. Its \term{dcterms:conformsTo} names the dated Recommendation of
the base's \acs{DCAT}, and its \term{owl:imports} names the release page for the % src: @5ffd56e B/compile/shacl.py:33-38
3.0.x bases and is absent for the others, so a validator loads the official shapes itself. The other artifacts % src: @5ffd56e B/data/base-profiles/*/profile.yaml:6 (homepage only for 3.0.0, 3.0.1); ISS K10, K11
are the minted vocabulary, \ac{LinkML} \cite{moxon2026linkml} for a visual profile editor\verify{the LinkML output
is valid LinkML (LinkML tooling not installed) and opens in the profile editor}, a \ac{DSV} profile \cite{klimek2026dsv} that, on a 3.0.1 % src: KCU §3.3-3.4; @5ffd56e B/compile/linkml_dsv.py:37-194
base, profiles the \acs{DSV} of \acs{DCAT-AP}~3.0.1 that Dataspecer generated \cite{klimek2024dataspecer}\verify{import of the DSV % src: @5ffd56e B/data/base-profiles/dcat-ap-3.0.1/dsv-index.json:2-3 (mff-uk DSV, generated by Dataspecer); KCU §3.4
output into Dataspecer}, the ProfileDelta, a trace matrix and documentation.

Eight rules then test what E4 asks: C1 legal refinement, C2 joint satisfiability, C3 % src: @5ffd56e B/conformance/check.py:5-21, 111-177, 250-266; CONF5:41-51
and C4 the minted and reused terms, C5 the trace, C6 \acs{SHACL} validity, which V6 already logs at compilation \cite{w3cshacl}, C7 a witness, and C8
the \acs{DCAT} state, failing on a deprecated or removed \acs{DCAT} term and warning on a later one. C7 builds an % src: @5ffd56e B/conformance/check.py:64-99
instance of each constrained class with a value for every required property and validates it with pySHACL % src: @5ffd56e B/conformance/check.py:181-247 (min count >= 1, :204-209; inference "none", :225); KCU §4.3
\cite{sommer2026pyshacl}, without inference,
against base and extension and against the base alone, warning when no valid witness results
(Sect.~\ref{sec:discussion}). The verdict is non-conformant when a rule fails at error level and conformant with % src: @5ffd56e B/conformance/check.py:264-266
warnings when only C7 or C8 warn.

The run stops at the review queue of waiting changes and conflicts: a person can accept,
reject, set back to undecided or edit any change and override any atom's disposition, and release requires an % src: @5ffd56e B/engine/pipeline.py:493-555
empty queue, a current conformance report that is not non-conformant and a recorded sign-off with a hash of the
included changes. An edit is a new version, verified and recorded as accepted; no later machine stage overwrites a % src: @5ffd56e B/engine/pipeline.py:611-635, 512-536
person's decision, and any later decision withdraws the sign-off. Each machine decision is shown with its % src: @5ffd56e B/engine/pipeline.py:115-126, 246-248, 283-285, 301-303, 637-638; ENG §4
evidence, against the over-trust in suggestions reported for OE-Assist \cite{lippolis2025oeassist}. % LIT G1-S16
\todo{tool: state or fix: unclear and unaddressed atoms never enter the queue and do not block the sign-off
(ENG §7.11; ISS E7.11; Q22)}

\subsection{Implementation and traceability}\label{sec:impl}

The pipeline runs in an open research workbench (Fig.~\ref{fig:workbench} in the appendix), where any stage can % src: @5ffd56e frontend/src/App.tsx:203-207, 241, 302; KCU §7
run alone and every intermediate result is on screen. An append-only SQLite blackboard records what every rule, model and
person did, each model call with its prompts, answer and tokens, and serves a replay cache; % src: @5ffd56e B/store/blackboard.py:1-11, 341-351; B/llm/client.py:1-12; KCU §5.3-5.4
the stages cooperate only through it, as the knowledge sources of Hearsay-II do % LIT G4-S01
\cite{erman1980hearsay,nii1986blackboard}, though in a fixed order. The trace matrix shows completeness, not correctness. The official % src: KCU §3.6
files of the nine \acs{DCAT-AP} releases with \acs{SHACL}, 1.1 to 3.0.1, and of \acs{DCAT}~2014, 2 and 3 are kept % src: DCS §1.1, §3; paper/inputs/tool-facts-main-5ffd56e.txt:6-14 (nine bases)
unmodified, with hashes\verify{1.1 and 1.2 record only the hash of the % src: DCS §1.1
zip file}; the 98 tests never call a model.\verify{98 tests pass at the frozen commit} % src: paper/inputs/tool-facts-main-5ffd56e.txt:24 (98 collected)

When the base release changes, a rebase moves the extension onto another release and its \acs{DCAT}: every change
is compared under both bases, human decisions stay, machine changes with a finding that asks for a decision are % src: @5ffd56e B/engine/rebase.py:73-180; B/engine/pipeline.py:557-609
flagged for review, and the extension is checked again, which withdraws the sign-off.
\todo{tool: the consolidation that follows a rebase re-decides flagged machine changes, and a lost covered-by-base
claim is reported but not re-proposed (ISS NEW-1, NEW-2; Q22)}
\todo{described commit: 5ffd56e, not yet in the paper branch (Q14)}

We took the design from 99 steps that we extracted from related work and fact-checked against their released % src: docs/literature-steps.md:3, 6
code, prompts, slides and specifications, some only against abstracts; 33 are implemented, 30 in part and 36 % LITD §0 caution 1 (abstract level for several G2 and G4 sources)
planned, and of the 50 from % src: docs/literature-steps.md:6 (33/30/36); G1 rows :12-36 (13 implemented, 9 in part), G2 rows :42-66 (9, 8): 50 steps, 22 and 17
ontology generation and requirements engineering, 22 are implemented and 17 in part (Table~\ref{tab:steps} in the % src: docs/literature-steps.md:12-36, 42-66 (G1 13+9, G2 9+8; recounted, outline 7.1)
appendix). The workbench is available at \todo{anonymized repository link}.
